\mainlanguage[de]

\setuplayout
  [backspace=2.0cm,
   height=middle,
   width=middle,
   leftmargin=1.8cm,
   footerdistance=0.0cm,
   rightmargindistance=0.2cm,
   rightmargin=1.8cm,
   bottom=0.5cm,
   leftmargindistance=0.2cm,
   header=2.1cm,
   headerdistance=0.5cm,
   footer=0.5cm,
   top=0cm,
   topspace=0.5cm,
   bottomspace=0.0cm]

\setupformulas[align=right]

\definehead[exercise][subsection]

\setuphead
  [exercise]
  [number=yes,
   style=,
   numberstyle=\ss\bf\tfb,
   textstyle=\ss,
   alternative=margin,
   numbercommand={\determineheadnumber[subsection]\currentheadnumber\gobbleoneargument}]

\setuphead
  [chapter]
  [number=no,
   style=,
   textstyle=\tfb\ss]

\setuphead
  [section]
  [number=no,
   style=,
   textstyle=\tf\ss]

\definemeasure[headerimageheight][\dimexpr(\headerheight-2ex)]

\definelayer[sabel]

\startsetups[sabel]

\setlayerframed[sabel][x=0mm,y=0mm][frame=off,offset=none]
{%
   \setupTABLE[start][style=\tfx]
   \setupTABLE[r][each][height=0.85\baselineskip]
   \setupTABLE[column][1][align=right, width=.4\textwidth, loffset=-2pt, toffset=-5pt]%
   \setupTABLE[column][2][align={center,lohi}, width=.2\textwidth]%
   \setupTABLE[column][3][align=left, width=.4\textwidth, roffset=-2pt, toffset=-5pt]%
   \bTABLE[frame=off]
   \bTR \bTD Sabel-GBS Fachakademie für Wirtschaft \eTD \bTD [nr=5] \dontleavehmode\externalfigure[../Logo/sabel.png][maxheight=\measure{headerimageheight}, frame=off]\dontleavehmode \eTD \bTD Wintersemester 2023/24 \eTD \eTR
   \bTR \bTD Mathematik Fachhochschulreife         \eTD                                                                                                                                     \bTD Klasse 4FA             \eTD \eTR
   \bTR \bTD B.~Donner                             \eTD                                                                                                                                     \bTD                        \eTD \eTR
   \bTR \bTD \strut                                \eTD                                                                                                                                     \bTD                        \eTD \eTR
   \bTR \bTD 1. Klausur                            \eTD                                                                                                                                     \bTD 13.12.2023             \eTD \eTR
   \eTABLE}

\stopsetups

\setupbackgrounds[header][text][bottomframe=on,setups=sabel,background=sabel]
\setuppagenumbering[location={footer,middle},style=\bfc]

\setupfootertexts
  [BE: 35 Arbeitszeit: 60-75min][Viel Erfolg!]
  [Viel Erfolg!][BE: 35 Arbeitszeit: 60-75min]


\definedescription[example][
  headstyle=\it, style=, align=left, location=top,
  width=broad, margin=1cm, indenting=never, text=\hskip-1cm]

\definemathmatrix[pmatrix][
  left={\left(\,}, right={\,\right)}]

%\usemodule[units]

\definemathcommand [kgV] [nolop] {\mfunction{kg\kern-.15em\hbox{V}}}

\setdigitmode 3

\setupitemgroup[itemize][1][n,repeat][width=1em]
\setupitemgroup[itemize][2][n,repeat][width=2em]
\setupitemgroup[itemize][3][n,repeat][width=3em]
\setupitemgroup[itemize][4][n,repeat][width=4em]

\startMPinclusions
  input ../include/mymacros.mp
\stopMPinclusions

\startreusableMPgraphic{graph1}
u := 1cm; % Unit length

path bbox;
xmin := -5; xmax := 5; ymin := -1.5; ymax := 4.5;
xdefmin := -4.9; xdefmax := 4.9;
xscale := 1; yscale := 1;
ybigticksmod := 2;
xbigticksmod := 2;

draw squaredblock(u*(xmax-xmin+1.5), u*(ymax-ymin+1.5), u*0.5) shifted (u*(xmin-0.5, ymin-0.5));

% Axes
pickup pencircle scaled 1bp;

drawoptions(withcolor black);
drawarrow (xmin*u, 0) -- (xmax*u, 0);
drawarrow (0, ymin*u) -- (0, ymax*u);
label.rt("\math{x}", (xmax*u, 0));
drawoptions(withcolor black);
label.rt("\math{y}", (0, ymax*u));

bbox = outerboundingbox currentpicture;

pickup pencircle scaled 2bp;

draw function(2, "x", "0.015625*((x+4)**2)*((x-4)**2)" , xdefmin,  xdefmax+0.1, 0.1)   scaled u withcolor black; 
%0.015625*(x+4)**2*(x-4)**2
draw ((-1.4641, 0) -- (-1.4641, 3) -- (1.4641, 3) -- (1.4641, 0)) scaled u;
label.lft("\math{f}", (-2.5u, 2u));
clip currentpicture to bbox;

pickup pencircle scaled 1bp;

labeloffset := (0.3+0.05)*u;

imax := trunc(yscale*ymax - ahlength/u);
for i = trunc(yscale*ymin) step 1 until imax :
  if i == 0 :
    label.llft("\math{" & decimal(i) & "}",   u*(0, i/yscale));
  elseif (i mod ybigticksmod) = 0 :
    draw ytick(0.3, (0, i/yscale)) scaled u;
    label.lft("\math{" & decimal(i) & "}",   u*(0, i/yscale));
  else:
    draw ytick(0.15, (0, i/yscale)) scaled u;
  fi;
endfor;


imax := trunc(xscale*xmax - ahlength/u);
for i = trunc(xscale*xmin) step 1 until imax :
  if i == 0 :
    label.llft("\math{" & decimal(i) & "}",   u*(0, i/yscale));
  elseif (i mod xbigticksmod) = 0 :
    draw xtick(0.3, (i/xscale, 0)) scaled u;
    label.bot("\math{" & decimal(i) & "}",   u*(i/xscale, 0));
  else:
    draw xtick(0.15, (i/xscale, 0)) scaled u;
  fi;
endfor;

\stopreusableMPgraphic

\starttext

Hinweis: Achten Sie auf eine übersichtliche und schöne äußere Form, korrekte mathematische Schreibweise und geben Sie bei Textaufgaben kurze Antworten an.


\startitemize

\item Von der Funktion \math{f(x) = \text{\digits{0.5}} x^3 - 5 x^2 + 16 x - 16} ist bekannt, dass sie die x-Achse an der Stelle 2 schneidet.
\startitemize
\item Bestimmen Sie die Nullstellen von \math{f} mit der entsprechenden Vielfachheit. (BE~5)
\item Geben Sie die Linearfaktorzerlegung(Faktorisierung) von \math{f} an. (BE~1)
\item Skizzieren Sie den den Verlauf des Graphen von \math{f} in ein Koordinatensystem mit korrekter Skalierung auf der \math{x}-Achse. (BE~2)
\stopitemize


\item Untersuchen Sie die Funktion \math{f(x) = x^4 - 3x^2}.
\startitemize
\item Weisen Sie die Symmetrieeigenschaften des Graphen nach. (BE~1)
\item Berechnen Sie die Nullstellen von \math{f(x)} und geben Sie deren Vielfachheit an. (BE~4)
\item Berechnen Sie die Koordinaten der Schnittpunkte des Graphen von \math{f} mit dem Graphen der Funktion \math{g(x) = 2x^2 - 4}. (BE~6)
\item Zeichnen Sie die Graphen von \math{f} und \math{g} im Bereich \math{\digits -2 \leq x \leq 2} in ein kartesisches Koordinatensystem unter Verwendung der bisherigen Ergebnisse. Achten Sie auf eine übersichtliche Zeichnung. (BE~5)
\stopitemize

\item Die Abbildung zeigt den obernen Teile des Giebels eines Barock-Hauses (Maße in \unit{m}) Der geschwungene Rand des Giebels wird durch eine ganzrationale Funktion \math{f} 4.Grades beschrieben.\\
\reuseMPgraphic{graph1}
\startitemize
\item Weisen Sie nach, dass \math{f} durch die Funktionsgleichung \math{\frac{1}{64}(x^4 - 32x^2 + 256)} gegeben ist. (BE~5)
\item Ein Fenster der Höhe \digits{3.0}\unit{m} soll in den Giebel eingepasst werden. Berechnen Sie exakt die maximal mögliche Breite des Fensters. (BE~6)

\stopitemize
\stopitemize

\stoptext
